%! Direct, Inverse \& Joint Variation — Unit 5, Lesson 5.7 — Student Packet Cover Sheet
\documentclass[10pt]{article}
\usepackage{algebra2-article}
\usepackage{algebra2-boxes}
\usepackage{ltablex}
\keepXColumns

\begin{document}

% Full-bleed forest banner
\begin{tikzpicture}[remember picture, overlay]
  \fill[forest] (current page.north west)
    rectangle ([yshift=-0.9in]current page.north east);
\end{tikzpicture}
\vspace{-0.8in}

\begin{center}
  {\color{white}\textbf{\LARGE Algebra 2: Shepherd}}\\[4pt]
  {\color{forestmid}\large\textbf{Unit 5 \quad Rational Functions}}\\[6pt]
  {\color{white}\textbf{\large Lesson 5.7 \quad Direct, Inverse \& Joint Variation}}
\end{center}

\vspace{0.22in}
\namedateperiod
\vspace{0.18in}

\begin{learningtargetbox}
\small
\begin{enumerate}[leftmargin=1.5em, itemsep=5pt]
  \item \ldots{} \textbf{decide} from a table whether two quantities are \textbf{directly
        proportional}, \textbf{inversely proportional}, or \textbf{neither} --- by testing whether the
        quotients $y\div x$ or the products $x\cdot y$ stay constant across \emph{every} row.
  \item \ldots{} \textbf{find} the constant of variation $k$, \textbf{write} the equation ($y=kx$ or
        $y=\frac kx$), and use it to answer the question that was asked.
  \item \ldots{} \textbf{recognize} each graph on sight: a direct variation is a line through the
        \emph{origin}; an inverse variation is Lesson 5.4's hyperbola, with the axes as its asymptotes.
  \item \ldots{} \textbf{interpret} $k$, a point, and the asymptotes in the words of the situation ---
        and say which answers the situation refuses even when the algebra allows them.
\end{enumerate}
\end{learningtargetbox}

\vspace{0.14in}

\begin{tocbox}
\small
\renewcommand{\arraystretch}{1.5}
\begin{tabularx}{\linewidth}{@{} c l X r @{}}
\rowcolor{forestbg}
\textbf{\#} & \textbf{Component} & \textbf{Description} & \textbf{Score} \\
\midrule
1 & Warm-Up        & Spiral review: quotient and product rows, the rational parent $\frac6x$, and one proportion & \blank{1.2cm} \\
2 & Guided Notes   & The two tests, direct and inverse worked in context, ``neither'' as a real answer, finding $k$, plus joint variation & \blank{1.2cm} \\
3 & Group Activity & \emph{Constant Quotient, Constant Product} --- table sorting, twin tables, a spring, rectangles of area 36, and an error analysis, by tier & \blank{1.2cm} \\
4 & Exit Ticket    & Quick check: name it and use it, correct a classmate's reasoning, plus an SOL-style item & \blank{1.2cm} \\
5 & Homework       & Tables, equations, two graphs to read, an error analysis, and the seesaw & \blank{1.2cm} \\
\midrule
  & \multicolumn{2}{r}{\textbf{Total}} & \blank{1.2cm} \\
\end{tabularx}
\end{tocbox}

\vspace{0.10in}

\begin{remindbox}
\small The unit ends where it began: with $y=\frac kx$. \textbf{Two questions decide every problem
today, and you answer them with arithmetic, not opinion.} Given a table, compute the quotients
$y\div x$ down the row and the products $x\cdot y$ down the row. \textbf{If the quotients are all the
same}, the quantities are \textbf{directly proportional}: $y=kx$, doubling $x$ doubles $y$, and the
graph is a line through the origin. \textbf{If the products are all the same}, they are
\textbf{inversely proportional}: $y=\frac kx$, doubling $x$ halves $y$, and the graph is the rational
parent from Lesson 5.4 with the axes as its asymptotes. \textbf{If neither row is constant, the answer
is ``neither''} --- and that is a complete, correct answer, not a sign you did something wrong. Two
warnings, because they cost more points than anything else on this lesson. First, \emph{``$y$ goes up
when $x$ goes up'' is not the test and ``$y$ goes down when $x$ goes up'' is not the test either.} A
phone bill of $\$15$ plus $\$5$ a gigabyte climbs perfectly steadily and is not a direct variation,
because at zero gigabytes you still owe $\$15$. Second, \emph{one row can never tell you anything} ---
every single pair of numbers has a quotient and a product, so the constant has to hold for the whole
table before you may name it. Finally, $k$ is never just a number. It has units and a meaning: dollars
per gallon, miles for the whole trip, pound-feet of balance. And when the answer has to describe
something real, some perfectly good solutions still get thrown out --- not because the algebra forbids
them, the way an extraneous solution is forbidden, but because the \textbf{situation} does.
\end{remindbox}

\end{document}
